\documentclass[10pt,a4paper]{amsart}
\usepackage[margin=19mm]{geometry}
\usepackage{amsmath,amssymb,mathtools}
\usepackage[scr=rsfs,cal=euler]{mathalfa}
\usepackage{xfrac}
\usepackage[unicode,hidelinks,bookmarks=false]{hyperref}
\newcommand{\ie}{i.e.\ }
\newcommand{\cf}{cf.\ }
\newcommand{\resp}{respectively\ }
\newcommand{\Subsection}{\subsection}
\providecommand{\nobreakdash}{\nobreak\hbox{-}}
\setlength{\emergencystretch}{2em}
\multlinegap0pt
\usepackage{bm}
\usepackage{bbm}
\usepackage{dsfont}
\usepackage{xspace}
\usepackage{cancel}
\usepackage{mathtools}
\usepackage{amsthm}
\makeatletter
\@ifundefined{theorem}{\newtheorem{theorem}{Theorem}}{}
\@ifundefined{lemma}{\newtheorem{lemma}[theorem]{Lemma}}{}
\makeatother
\title[Residual stratification and the Cantor-Bendixson structures ]{Residual stratification and the Cantor-Bendixson structures of dual algebraic coframes}
\author{Silv\`ere Gangloff, Alonso Nu\~nez}
\date{Source: arXiv:2605.04851v3 (CC BY 4.0)}
\begin{document}
\maketitle

\noindent\emph{Source and licence.} Residual stratification and the Cantor-Bendixson structures of dual algebraic coframes. Version arXiv:2605.04851v3 (CC BY 4.0), distributed under the Creative Commons Attribution 4.0 International licence, as recorded in the arXiv licence metadata for this exact version and shown on the abstract page https://arxiv.org/abs/2605.04851v3. Full rights notice: licenses/arxiv-2605.04851-CC-BY-4.0.txt.\par\smallskip

\noindent\textit{Editorial scope.} Counted unit: the Lemma whose source label is \texttt{lemma.subelement.decomposition} in the article (source0.tex lines 1085--1088, source label \texttt{lemma.subelement.decomposition}), delivered together with the article's own complete original proof of it (source0.tex lines 1090--1187, one \texttt{proof} environment of 98 lines). No mathematics is written, repaired or re-derived: statement and proof are the article's own text line for line. Required context retained uncounted: lemma.minmax (lines 1068--1071). Every definition, display and label of those lines is retained unchanged. Omitted from this package and not redistributed: 1--1067, 1072--1084, 1089--1089, 1188--2078. Nothing inside a retained excerpt is omitted or altered: every retained body is delivered exactly as the article's lines are written, so no omission receipt of any kind accompanies this package. Citations: the retained excerpts carry no citation command at all, so no bibliography entry of the article is transcribed and no reference is redistributed. Reference adaptation: none was needed and none was made; every reference inside the delivered unit resolves inside it, so no unresolved reference or citation is produced. Printed slips kept literal: no printed word, symbol, hypothesis or number of the counted statement or of its proof was altered. Layout and class: the article's own class is \texttt{article} with packages including inputenc, authblk, amsthm, xcolor, ulem, latexsym, amscd, amsmath, amssymb, amsfonts, tikz, stmaryrd, hyperref, tcolorbox; the wrapper is the lane's pinned \texttt{amsart} editorial wrapper at 10pt on A4, which restates the theorem-like environments the retained text needs and adds the article's own macro definitions that the retained text needs (none), with extra wrapper packages (bm, bbm, dsfont, xspace, cancel, mathtools). The delivered numbering is the unit's own. Assets: no retained span contains a figure environment, a graphics-inclusion command, a PDF-page inclusion, a table or a TikZ picture; no image file is copied into this package, the package declares no asset dependency and no image file is redistributed with it. Licence: version arXiv:2605.04851v3 of this article is distributed under the Creative Commons Attribution 4.0 International licence, as recorded in the arXiv licence metadata for this exact version and shown on the abstract page https://arxiv.org/abs/2605.04851v3. A full rights notice is delivered at licenses/arxiv-2605.04851-CC-BY-4.0.txt. Characters: every retained excerpt is pure ASCII, so the delivered engine, XeLaTeX with its default Latin Modern text font, reports no missing character.
\begin{lemma}\label{lemma.minmax}
Let \(S\) be any set. Let \((x_s)_{s \in S}\) and \((z_s)_{s \in S}\) be families in \(L\). We have 
\[\bigwedge_{s \in S} (x_s \vee z_s) \leq(\bigvee_{s \in S} x_s)\vee(\bigwedge_{s \in S} z_s).\]
\end{lemma}

\begin{lemma}\label{lemma.subelement.decomposition}
    For every $x \in L$ and $ z \in \mathop{\downarrow}x$, we have 
    \[z = (z \wedge \texttt{c}(x)) \vee \bigvee_{\underset{s \le z}{s \in \delta(x)}} s.\]
\end{lemma}

\begin{proof}
    \textbf{1.} 
    % We first prove that when $z \ge \texttt{c}(x)$, we have  \[z = \texttt{c}(x) \vee \bigvee_{\underset{s \le z}{s \in \delta(x)}} s.\]
    {We first prove the statement when \(z\ge \texttt{c}(x)\). For every \(\alpha<\texttt r(x)\), set

\[u_{\alpha} := x^{(\alpha)}\vee \left(\bigvee_{\beta < \alpha} \bigvee_{\underset{s \le z}{s \in \texttt{s}_{\beta}(x)}}s\right)\]

We prove by transfinite induction on \(\alpha\) that \(z\leq u_\alpha\). We have $u_0 = x$, therefore $z \le u_{0}$ is straightforward. Suppose now that the statement holds for every \(\beta<\alpha\). Assume that \(\alpha\) is a limit ordinal. We apply Lemma \ref{lemma.minmax}, where $S = \texttt{r}(x)$, and for all $\beta \in \texttt{r}(x)$,
\[x_{\beta} := \bigvee_{\lambda < \alpha} \bigvee_{\underset{s \le z}{s \in \texttt{s}_{\lambda}(x)}}s \qquad \text{and} \qquad z_{\beta} := x^{(\beta)}.\]
We have $\bigvee_{\beta < \alpha} x_{\beta} = x_{\alpha}$ and 
$\bigwedge_{\beta < \alpha} z_{\beta} = z_{\alpha}$.
We thus have \(u_\alpha \ge \bigwedge_{\beta<\alpha}u_\beta\). Since \(z\leq u_\beta\) for every \(\beta<\alpha\), we obtain \(z\le u_\alpha\).} 
%     For every $\alpha < \texttt{r}(x)$, set: 
%     \[z_{\alpha} := \texttt{c}(x) \vee \mu(x^{(\alpha)}) \vee \left(\bigvee_{\beta < \alpha} \bigvee_{\underset{s \le z}{s \in \texttt{s}_{\beta}(x)}}s\right) \vee \left( \bigvee_{{s \in \texttt{s}_{\alpha}(x)}}s\right)\] 
%     {By Lemma \ref{lemma.minmax},} it is sufficient to see that for every ordinal $\alpha < \texttt{r}(x)$, we have $z \wedge z_{\alpha}= z_{\alpha}$.
%     We prove this by transfinite induction. Consider  $\alpha < \texttt{r}(x)$ such that for all $\beta < \alpha$, we have $z \wedge z_{\beta} =  z_{\beta}$, then $z \le z_{\beta}$. If $\alpha$ is a limit ordinal, 
% %by taking limits, since 
%     %$\le$ is closed, 
%     {by Lemma \ref{lemma.minmax}} we obtain
%     that $z \le z_{\alpha}$, which means that $z \wedge z_{\alpha} = z$. 
Otherwise we set $\lambda$ such that $\lambda + 1 = \alpha$. By distributivity, for every finite subset {$P \subset \texttt{s}_{\lambda}(x)$}:
%$P \subset \texttt{s}_{\alpha}(x)$: 
% \[
% \begin{aligned}
% z \wedge z_{\alpha}
%   ={}&(z \wedge \texttt{c}(x))
%    \vee (z \wedge \mu(x^{(\alpha)}))
%    \vee \left(z \wedge \bigvee_{\beta < \alpha}
%         \bigvee_{\substack{s \in \texttt{s}_{\beta}(x)\\ s \le z}} s\right) \\
%    &\vee \left(\bigvee_{s \in P} z \wedge s\right)
%    \vee \left(z \wedge \bigvee_{s \in \texttt{s}_{\alpha}(x)\setminus P} s\right).
% \end{aligned}
% \]
{
\[
\begin{aligned}
z = z \wedge u_{\lambda}
  ={}&(z \wedge \mu(x^{(\lambda)}))
   \vee \left(z \wedge \bigvee_{\beta < \lambda}
        \bigvee_{\substack{s \in \texttt{s}_{\beta}(x)\\ s \le z}} s\right) \\
   &\vee \left(\bigvee_{\underset{s \le z}{s \in P}} z \wedge s\right) \vee \left(z \wedge\bigvee_{\underset{s \not \le z}{s \in P}} s\right)
   \vee \left(z \wedge \bigvee_{s \in \texttt{s}_{\lambda}(x)\setminus P} s\right).
\end{aligned}
\]

For $s \in P$, if $z \wedge s < s$ then $z \wedge s \le \mu(x^{(\lambda)})$, and thus 
$z \wedge s \le z \wedge \mu(x^{(\lambda)})$. We can thus rewrite: 

\[
\begin{aligned}
z \wedge u_{\lambda}
  ={}&(z \wedge \mu(x^{(\lambda)}))
   \vee \left(z \wedge \bigvee_{\beta < \lambda}
        \bigvee_{\substack{s \in \texttt{s}_{\beta}(x)\\ s \le z}} s\right) \\
   &\vee \left(\bigvee_{\underset{s \le z}{s \in P}} z \wedge s\right) \vee \left(z \wedge \bigvee_{s \in \texttt{s}_{\lambda}(x)\setminus P} s\right).
\end{aligned}
\]

We apply again Lemma \ref{lemma.minmax}, with $S$ equal to the set of finite subsets of $\texttt{s}_{\lambda}(x)$, and for all $P \in S$, 
\[x_P := (z \wedge \mu(x^{(\lambda)}))
   \vee \left(z \wedge \bigvee_{\beta < \lambda}
        \bigvee_{\substack{s \in \texttt{s}_{\beta}(x)\\ s \le z}} s\right) \vee \left(\bigvee_{\underset{s \le z}{s \in P}} z \wedge s\right),\]
and 
\[z_P := \left(z \wedge \bigvee_{s \in \texttt{s}_{\lambda}(x)\setminus P} s\right)\]

By application of the lemma, 
we have $z \wedge u_{\lambda} \le \left(\bigvee_{P \in S} x_P\right) \vee \left(\bigwedge_{P \in S} z_P\right)$. We have 
\[\bigwedge_{P \in S} z_P \le z \wedge \bigwedge_{\underset{P finite}{P \subset \mathcal{M}(x^{(\lambda)})}} \bigwedge_{m \in \mathcal{M}(x^{(\lambda)})\setminus P} m \le z \wedge \mu(x^{(\lambda)})\]
We thus have also 
\[z \le z \wedge \left(\mu(x^{(\lambda)}) \vee  \left(\bigvee_{\beta \le \lambda}
        \bigvee_{\substack{s \in \texttt{s}_{\beta}(x)\\ s \le z}} s\right)\right) = z \wedge u_{\alpha}.\]
In particular, $z \le u_{\alpha}$.
}
%For $s \in P$, if $z \wedge s < s$ then $z \wedge s \le \mu(x^{(\alpha)})$, and thus 
%$z \wedge s \le z \wedge \mu(x^{(\alpha)})$.
% Since we have 
% \[\bigwedge_{\underset{P finite}{P \subset \texttt{s}_{\alpha}(x)}} \bigvee_{s \in \texttt{s}_{\alpha}(x)\cap P} \le \bigwedge_{\underset{P finite}{P \subset \mathcal{M}(x^{(\alpha)})}} \bigwedge_{m \in \mathcal{M}(x^{(\alpha)})\setminus P} m \le \mu(x^{(\alpha)}),\]

% Using Lemma \ref{lemma.formula.maximal} on $\mu(x^{(\alpha)})$, we get
% $z \wedge z_{\alpha} = z \wedge z_{\alpha+1} = z$. 

{
\textbf{2.} We apply the first point on $z \vee \texttt{c}(x)$, and get: 
\[z \vee \texttt{c}(x) = \texttt{c}(x) \vee \bigvee_{\underset{s \le z}{s \in \delta(x)}} s.\]
Then by distributivity 
\[z = z \wedge (z \vee \texttt{c}(x)) = (z \wedge \texttt{c}(x)) \vee \left(z \wedge  \bigvee_{\underset{s \le z}{s \in \delta(x)}} s\right).\]

Since the last term is below $z$, 
we get 
\[z = (z \wedge \texttt{c}(x)) \vee \bigvee_{\underset{s \le z}{s \in \delta(x)}} s.\]
% For an arbitrary \(z\leq x\), set
% \[ z':=(z\wedge\texttt c(x))\vee\bigvee_{\substack{s\in\delta(x)\\ s\leq z}}s.\]
% Clearly \(z'\leq z\). Moreover,
% \[z\vee\texttt c(x)=\texttt c(x)\vee\bigvee_{\substack{s\in\delta(x)\\ s\leq z}}s=z',\]
% and taking the meet with \(z\) gives \(z=z\wedge(z\vee\texttt c(x))=z\wedge z'=z'\). Therefore
% \[z=(z\wedge\texttt c(x))\vee\bigvee_{\substack{s\in\delta(x)\\ s\leq z}}s,\text{ as required.}\]
}
\end{proof}

\end{document}
